%---------------------------------------------------------------------------
%
%                          Vorlage der Arbeitsgruppe
%             Computer Vision and Pattern Recognition Group (CVPR)
%                           der Universität Münster
%                         http://cvpr.uni-muenster.de
%
%---------------------------------------------------------------------------
% Geeignet für:
%  - Seminararbeiten
%  - Bachelorarbeiten
%  - Masterarbeiten
%---------------------------------------------------------------------------
% Autoren:
%  - Daniel Tenbrinck
%  - Fabian Gigengack
%  - Michael Schmeing
%  - Lucas Franek
%  - Andreas Nienkötter
%---------------------------------------------------------------------------
% Version:
%  - 1.0.3 (05.10.2016)
%	 - Ersetzung von veralteten Befehlen durch Aktuelle
%	 - Einige ausführlichere Beispiele
%    - Einführung von listings
%    - Aktuelle Version der Eidesstattlichen Erklärung
%  - 1.0.2 (09.09.2011)
%    - Titelblatt um Matrikelnummer und Studiengang ergänzt
%  - 1.0.1 (05.07.2011)
%---------------------------------------------------------------------------
% 
% "THE BEER-WARE LICENSE" (Revision 42):
% The above mentioned authors wrote this file. As long as you retain this
% notice you can do whatever you want with this stuff. If we meet some day,
% and you think this stuff is worth it, you can buy us a beer in return.
%  --------------------------------------------------------------------------

\documentclass[a4paper, twoside, 12pt, ngerman]{scrbook} % Layout-Einstellungen für das Dokument

\usepackage[utf8]{inputenc} % UTF-8 Codierung
\usepackage[ngerman]{babel} % Deutsche Beschriftung

\usepackage{graphicx} % Um Bilder einzufügen
\usepackage{subfigure} % Um mehrere Bilder in eine figure einzufügen
\usepackage{amssymb, amsmath} % Für Mengensymbole und über Gleichheitszeichen schreiben
\usepackage{verbatim} % Um Quellcode in das Dokument einzufügen.
\usepackage{xcolor} % Für Farben
\usepackage[linkbordercolor=blue]{hyperref} % Für Links im Dokument
\usepackage{algorithmic} % Für Pseudo-Code
\usepackage{algorithm} % Wrapper für Pseudo-Code
\usepackage[font={small}, labelfont=bf]{caption} % kleine Bildunterschriften
\usepackage{geometry} % Für Feinanpassungen des Layouts

\usepackage{listings} % Für Code-Listings
\renewcommand{\lstlistingname}{Quelltext} %Ändert die Überschrift von Listing nach Quelltext

% Einstellungen für Abstand an den Rändern
\geometry{a4paper,left=35mm,right=35mm,top=20mm,bottom=20mm, includeheadfoot}

\begin{document}
\pagenumbering{roman}

% Titelblatt
\begin{titlepage}
\begin{centering}
%\vspace*{1cm}
\includegraphics[width=12cm]{./img/wwu-logo-neu.pdf}

\vspace{2cm} 

{\LARGE
	\textbf{Ein graphtheoretischer Ansatz für das \emph{multiple sequence Alignment}-Problem}\\[1.2cm]
}

{\large
	Bachelorarbeit\\[2cm]
}

{\large
	vorgelegt von:
}

{ \Large
	\textbf{Joschka Strüber}\\[1cm]
}

{\large
	Matrikelnummer: 418702\\[2mm]
}

{\large
	Studiengang: B.Sc. Informatik\\[1cm]
}
    
{\large
	Thema gestellt von:
}

{\Large
	\textbf{Prof. Dr. Jan Vahrenhold}\\[1cm]
}
                               
{\large
	Arbeit betreut durch:
}

{\Large
	\textbf{Prof. Dr. Jan Vahrenhold \\ Prof. Dr. Xiaoyi Jiang}\\[1cm]
}

{\large
Münster, \today
}
\vfill
\end{centering}
\end{titlepage}

% Inhaltsverzeichnis
\tableofcontents

\cleardoublepage
\pagenumbering{arabic}

% Die Hauptkapitel der Arbeit
\chapter{Einleitung}
\label{ch:einleitung}

\section{Multiple Sequence Alignments}

\section{Einsatzgebiete}

\section{Komplexität}

\chapter{Methodik}
\label{ch:methodik}

Im Methodik-Kapitel werden die mathematischen Ausführungen des Verfahrens bzw. des Algorithmus vorgestellt, jedoch technische Details zur Umsetzung und Implementierung (falls nötig) auf ein darauffolgendes Kapitel verlagert.


\section{Mathematische Notation}
\label{s:notation}

Mathematische Formeln können mittels der \verb+\begin{align}...\end{align}+ Umgebung gesetzt werden:

\begin{align}
f(n) & =
	\begin{cases}
		n/2, & \text{wenn }n\text{ gerade,}\\
		3n+1, & \text{wenn }n\text{ ungerade.}
	\end{cases}
\label{eq:f} \\
%
g(n) & = \frac{n}{2} \label{eq:g}
\end{align}

\section{Algorithmus}
\label{s:algorithmus}

Eigene Algorithmen beschreibt man am Besten mit Hilfe von Pseudo-Code und dem Paket \verb+algorithm+.

\begin{algorithm}
\caption{Algorithmus}
\label{alg:alg}
\begin{algorithmic}
\algsetup{indent=2em}

\REQUIRE Argument $n\in\mathbb{N}$
\STATE $a = 0$
\FOR{ $i=0,\dots,n$}
	\STATE $a = a + 1$
\ENDFOR
\RETURN $a$
\end{algorithmic}
\end{algorithm}

\chapter{Ergebnisse}
\label{ch:ergebnisse}

Dieses Kapitel sollte die Ergebnisse beinhalten, die mit den Methoden aus \autoref{ch:methodik} erstellt wurden.

\begin{table}[h]
\caption{Beispieltabelle}
\begin{center}
	\begin{tabular}{|c||c|c|}
		\hline
		Spalte1 & Spalte2 & Spalte3 \\ 
		\hline\hline
		   1    &    2    &    3    \\ 
		\hline
	\end{tabular}
\end{center}
\label{tbl:table}
\end{table}

\chapter{Diskussion}
\label{ch:diskussion}

In diesem Kapitel werden die zuvor vorgestellten Ergebnisse der Arbeit diskutiert. Häufig wird es mit dem Fazit zusammengelegt.

\chapter{Fazit}
\label{ch:fazit}

\section{Zusammenfassung}

\section{Future Works} \label{sec:fut_work}

\begin{itemize}
	\item Statt paarweisen Alignments, Alignments von je drei Sequenzen berechnen
	\item Conditional Random Fields statt Gewichtsfunktionen
	\item Bessere Heuristik zum Löschen der Kanten aus dem Sukzessorgraph benutzen (gerade wenn Tests zeigen, dass oft sehr viele Kanten gelöscht werden) - das könnte vor allem dann vorkommen, wenn die Sequenzen große Überkreuzungen enthalten, weil man dann große Zyklen mit sehr hohen Kantengewichten hat. -> mögliches Paper?
	\item Statt DIALIGN 2 DIALIGN TX zwischen den Ankerpunkten nutzen
	\item Wir wollen aus dem Sukzessorgraph möglichst wenige Knoten löschen. Dafür sind Sites, die aber als einzige in ihren Knoten vorkommen, irrelevant. 
		\begin{enumerate}
			\item Jede Kante hinter Knoten hat Gewicht von Anzahl an Sites - 1 => bevorzugt Zuweisungen über möglichst viele Sequenzen hinweg. Es kann aber sein, dass dadurch viele kleine Zuweisungen aufgelöst werden
			\item Jede Kante hinter Knoten hat Gewicht 1, falls mehr als eine Sequenz beteiligt und Gewicht 0, falls nicht => bevorzugt viele kleine Alignments und minimiert Anzahl der Löschungen
		\end{enumerate}
	\item Teile des Algorithmus lassen sich gut parallelisieren:
		\begin{enumerate}
			\item Die $(\frac{1}{2}*(n^2-n))$-vielen paarweisen Alignments lassen sich parallelisiert berechnen.
			\item Die Überlappgewichte lassen sich gut parallelisiert berechnen, weil keins der Ergebnisse von denen der anderen abhängt. Paralleles Lesen der paarweisen Alignments ist kein Problem.
			\item Min-Cuts können innerhalb jeder Zusammenhangskomponente parallelisiert berechnet werden.
			\item Bei der alten Methode können auch die kürzesten Pfade durch den Sukzessionsgraphen parallelisiert werden, bei den beiden neuen Ansätzen nicht. Das ist aber nicht so schlimm, weil dieser Abschnitt die Laufzeit nicht dominiert.
		\end{enumerate}
	\item Ersetzen der Felder für die \emph{Waisen} durch eine bessere Datenstruktur
	\item Hinzufügen von Kanten im \emph{Alignmentgraphen} für ganze \emph{Fragmente}
\end{itemize}

% Literaturverzeichnis
\bibliographystyle{unsrtdin}
\bibliography{Quellen}

% Eidesstattliche Erklärung
\chapter*{Eidesstattliche Erklärung}

% Die Aktuelle Version der Eidesstättlichen Erklärung kann beim zuständigen Prüfungsamts gefunden werden.
% Nachfolgend ist die Version für Bachelor und Master des Prüfungsamts Mathe/Informatik vom 05.10.2016

Hiermit versichere ich, dass die vorliegende Arbeit über \textit{\glqq Ein graphtheoretischer Ansatz für das multiple sequence alignment-Problem\grqq} selbstständig verfasst worden ist, dass keine anderen Quellen und Hilfsmittel als die angegebenen benutzt worden sind und dass die Stellen der Arbeit, die anderen Werken – auch elektronischen Medien – dem Wortlaut oder Sinn nach entnommen wurden, auf jeden Fall unter Angabe der Quelle als Entlehnung kenntlich gemacht worden sind.

\vspace{1cm}

\parbox{20em}{\hrulefill}

Vorname Nachname, Münster, \today

\vspace{1cm}

Ich erkläre mich mit einem Abgleich der Arbeit mit anderen Texten zwecks Auffindung von Übereinstimmungen sowie mit einer zu diesem Zweck vorzunehmenden Speicherung der Arbeit in eine Datenbank einverstanden.

\vspace{1cm}

\parbox{20em}{\hrulefill}

Vorname Nachname, Münster, \today

\end{document}